\documentclass[11pt,oneside]{report}

\usepackage[utf8]{inputenc}
\usepackage[T1]{fontenc}
\usepackage[margin=1in]{geometry}
\usepackage{array}
\usepackage{longtable}
\usepackage{amsmath}
\usepackage{amssymb}
\usepackage{hyperref}

\hypersetup{
  colorlinks=true,
  linkcolor=black,
  citecolor=black,
  urlcolor=blue
}

\setlength{\parindent}{0pt}
\setlength{\parskip}{6pt}
\sloppy
\emergencystretch=3em
\renewcommand{\arraystretch}{1.2}
\newcolumntype{L}[1]{>{\raggedright\arraybackslash}p{#1}}
\newcommand{\project}{MemSieve}
\newcommand{\missing}{\textbf{Missing / Must Decide: }}
\newcommand{\evidence}{\textbf{Evidence required: }}
\newcommand{\risk}{\textbf{Risk: }}
\newcommand{\fix}{\textbf{Required fix: }}

\title{\textbf{\project: Implementation and Paper Roadmap}\\
\large A Harsh Research Audit, System Blueprint, Experimental Plan, and Submission Manual}
\author{Prepared from \texttt{paper\_roadmap.pdf}}
\date{May 2, 2026}

\begin{document}
\maketitle
\tableofcontents

\chapter*{Source Discipline and Audit Position}
\addcontentsline{toc}{chapter}{Source Discipline and Audit Position}

This document is a standalone implementation-and-paper manual derived from the repository file
\texttt{paper\_roadmap.pdf}. The repository also contains \texttt{AeroMind Paper.pdf}, but this
manual does not use it because the task explicitly names \texttt{paper\_roadmap.pdf} as the source
report. Any later use of other private or unpublished material must be declared and reviewed for
overlap risk before submission.

The source report is strong in its research framing: it correctly identifies that the publishable
unit is not ``memory can be attacked,'' but a defense-first runtime architecture for mutable LLM
agent memory. It also makes the right strategic call: the paper should focus on a compact,
cross-framework, cross-subsystem memory firewall with scalability evidence, not a broad benchmark,
new attack taxonomy, or heavyweight formal system.

The source report is weak as an implementation document. It does not specify exact APIs,
thresholds, task schemas, framework integration contracts, datasets, statistical tests,
reproducibility artifacts, or acceptance criteria for individual modules. It cites many public
threads of work but does not provide verified bibliographic metadata. It proposes a two-week
execution schedule, which is plausible only for a disciplined MVP and implausible for a broad
main-conference paper with full baseline reproduction.

\textbf{Operating rule for this manual:} when the report does not specify something, this document
marks it as \missing{} rather than pretending the detail exists. Implementation decisions below are
derived from the report's thesis and converted into concrete contracts so a developer can begin
coding immediately.

\chapter{Research Vision and Paper Thesis}

\section{Executive Research Vision}

\textbf{Core research idea.} \project{} is a provenance-aware memory firewall for LLM agents. It sits
between agent frameworks and their memory APIs, forcing every memory write, retrieval, context
assembly, tool action, and cross-agent memory propagation event through a common defensive pipeline.
The core architectural insight is to treat agent memory as a mutable write--retrieve--propagate
system rather than as a passive vector index or prompt suffix.

\textbf{Problem solved.} Modern LLM agents increasingly store durable state: user facts, project
preferences, retrieved passages, prior trajectories, shared team notes, tool outputs, summaries, and
coordination artifacts. These memory records can be contaminated by poisoned corpora, indirect
prompt injection, low-trust user content, compromised peer agents, or self-reinforcing summaries.
Current defenses are fragmented: they usually protect one stage, one substrate, or one framework.
\project{} aims to reduce contaminated writes, unsafe retrieval, contradiction-induced failures,
tool-use abuse conditioned on memory, and cross-agent propagation.

\textbf{Why the problem matters now.} The report argues that public agent frameworks now expose
durable memory as a first-class runtime feature. This changes the security problem from prompt
safety at one interaction to persistent state safety across sessions and agents. The important
shift is persistence: a contaminated memory record can keep influencing future reasoning after the
original attacker interaction has disappeared.

\textbf{Expected paper-level contribution.} The expected contribution is a compact systems/security
paper that demonstrates a reusable runtime layer over public memory APIs. The paper should show that
one canonical schema plus multi-stage gates can materially lower unsafe retrieval and propagation
across vector/RAG memory, semantic fact memory, episodic/example memory, and shared coordination
memory, while preserving clean utility and keeping overhead bounded as memory size and agent count
grow.

\textbf{Why this is publishable rather than merely engineering.} The publishable claim is not that a
developer can wrap memory APIs. The publishable claim is that memory defense needs a new evaluation
object: the full write--retrieve--propagate lifecycle across heterogeneous memory subsystems. A
paper becomes credible if it shows claim-evidence alignment: generic metrics, real framework
adapters, cross-subsystem scenarios, ablations, and scalability curves. Without that evidence,
\project{} is only a well-named middleware prototype.

\section{Paper Thesis}

The strongest thesis is:

\begin{quote}
LLM agent memory should be defended as a stateful runtime subsystem. A framework-agnostic middleware
that normalizes memory records into a provenance- and lineage-aware schema, gates writes, reranks
retrievals under risk, renders retrieved content as typed evidence, verifies memory-conditioned tool
actions, and controls cross-agent propagation can substantially reduce memory contamination while
preserving clean utility and scaling acceptably.
\end{quote}

The paper must not claim general immunity. It should claim measurable reduction under a bounded
threat model.

\section{Non-Negotiable Scope}

\begin{itemize}
  \item MVP frameworks: LangGraph, LlamaIndex, and CrewAI.
  \item Optional framework: AutoGen only if integration is low-friction.
  \item Stretch framework: Letta/MemGPT only if it is practically zero-friction.
  \item MVP memory subsystems: vector/RAG, semantic fact, episodic/example, shared coordination.
  \item MVP defense stages: adapter, write gate, quarantine/audit, retrieval gate, context shaping,
        pre-tool verification, propagation controller.
  \item MVP evidence: clean utility table, attack reduction table, ablation table, propagation
        result, scalability curves, reproducibility package.
\end{itemize}

\chapter{Novelty, Gaps, and Contributions}

\section{Strongest Novelty Claims Extracted from the Report}

\begin{longtable}{|L{0.25\textwidth}|L{0.35\textwidth}|L{0.30\textwidth}|}
\hline
\textbf{Claim} & \textbf{Why it is potentially strong} & \textbf{Evidence needed} \\
\hline
Cross-lifecycle memory defense & Most defenses target retrieval, prompts, or tools separately.
The report's strongest move is to defend writes, retrievals, context assembly, tool actions, and
propagation as one lifecycle. & Ablations showing that no single stage explains all gains. \\
\hline
Cross-subsystem coverage & Vector-only defenses are easy to dismiss. Covering semantic,
episodic, and shared coordination memory makes the paper meaningfully about agent memory. & Main
results table must include non-RAG rows, not appendix-only evidence. \\
\hline
Framework-agnostic runtime layer & A common schema plus adapters can make the architecture
portable across public frameworks. & Working adapters for at least three frameworks with identical
scenario templates and logging. \\
\hline
Propagation-aware memory security & Shared memory and peer-agent summaries make contamination
spread measurable. & Lineage graph, propagation depth, fan-out, and agent-count scaling. \\
\hline
Scalability-first evaluation & Security that only works on tiny stores is not deployable. & Curves
over memory size, poison ratio, top-k, and agent count. \\
\hline
\end{longtable}

\section{Weak, Risky, or Already Covered Claims}

\begin{longtable}{|L{0.23\textwidth}|L{0.34\textwidth}|L{0.33\textwidth}|}
\hline
\textbf{Risky claim} & \textbf{Why it is weak} & \textbf{Rewrite or remove} \\
\hline
``We introduce memory attacks.'' & The report itself says public work already covers indirect
prompt injection, RAG poisoning, memory injection, and propagation. & Remove. Position attacks as
evaluation scenarios grounded in public threat categories. \\
\hline
``Framework-agnostic defense.'' & This is only true if adapters actually run. Otherwise reviewers
will treat it as architecture rhetoric. & Claim ``portable across three public frameworks via thin
adapters'' after demonstrating LangGraph, LlamaIndex, and CrewAI. \\
\hline
``Defends all memory.'' & The report explicitly recommends four stable subsystems and says
procedural/scratchpad memory is an extension. & Claim coverage of four practical memory subsystems;
mark procedural memory as future work. \\
\hline
``Novel detector.'' & The report warns against a monolithic learned detector. A detector-focused
claim would be a schedule and review trap. & Claim interpretable, calibratable gates and optional
LLM checks, not a new classifier. \\
\hline
``Strong conference main-track readiness in two weeks.'' & The report is blunt that this is
unlikely unless most code already exists. & Target a compact systems/security paper or workshop
submission first; build a main-track version only after stronger evaluation. \\
\hline
\end{longtable}

\section{Publishable Novelty Statement}

\project{} is a framework-portable memory firewall that defends LLM agent memory across the
write, retrieve, reason, act, and propagate lifecycle. It normalizes heterogeneous memory records
into a common schema carrying provenance, trust priors, contradiction signals, and lineage edges;
quarantines or down-trusts suspicious writes; reranks retrieved candidates with risk-aware scoring
and source diversity; renders memory as typed evidence rather than executable instruction; verifies
memory-conditioned tool actions; and limits cross-agent spread through lineage-aware propagation
control. The novelty is the combination of lifecycle coverage, cross-subsystem evaluation, public
framework adapters, generic memory-security metrics, and scalability evidence.

\section{Concrete Paper Contributions}

\begin{longtable}{|L{0.06\textwidth}|L{0.29\textwidth}|L{0.29\textwidth}|L{0.26\textwidth}|}
\hline
\textbf{\#} & \textbf{Contribution} & \textbf{What the paper must prove} & \textbf{Primary evidence} \\
\hline
1 & A lifecycle defense architecture for mutable agent memory. & Multi-stage defense reduces
contaminated writes, unsafe retrieval, unsafe context, unsafe tool action, and propagation more than
single-stage baselines. & Main attack table plus ablation table. \\
\hline
2 & A canonical provenance and lineage schema for heterogeneous agent memory. & One schema can
represent vector/RAG, semantic, episodic, and shared coordination records across three frameworks. &
Schema definition, adapter tests, cross-framework logs. \\
\hline
3 & A cross-subsystem memory-security evaluation harness. & The same scenarios and metrics run
across multiple frameworks and memory types without bespoke hand-tuning per row. & Experiment runner,
scenario manifests, reproducibility artifacts. \\
\hline
4 & New evaluation metrics for memory contamination and propagation. & Unsafe retrieval alone is
insufficient; contaminated context and propagation depth capture stronger failure modes. & Metric
definitions, examples, and plots linking metrics to downstream outcomes. \\
\hline
5 & Scalability evidence for security under memory growth and team growth. & Security gains remain
meaningful as memory size, poison ratio, top-k, and agent count increase, with acceptable overhead. &
Scalability curves and security-cost frontier. \\
\hline
\end{longtable}

\section{Gaps That Must Be Filled Before Coding}

\begin{itemize}
  \item \missing{} Exact model provider, model version, embedding model, and cost budget.
  \item \missing{} Exact framework versions and whether adapters target Python or JavaScript APIs.
  \item \missing{} Threshold values for write admission, retrieval risk, quarantine, trust decay, and
        confirmation.
  \item \missing{} Public corpus choice or synthetic corpus generator parameters.
  \item \missing{} Tool-action risk taxonomy and which stub tools count as high impact.
  \item \missing{} Statistical test policy: confidence intervals, paired tests, bootstrap, or only
        descriptive statistics.
  \item \missing{} Submission venue, page limit, artifact expectations, and anonymity policy.
\end{itemize}

\chapter{System Architecture}

\section{Architecture Overview}

\project{} sits between an agent runtime and its native memory implementation. The framework should
believe it is calling a normal memory API; \project{} intercepts and records every memory event.

\begin{figure}[h]
\centering
\fbox{\begin{minipage}{0.92\textwidth}
\textbf{Figure 1 textual architecture placeholder.}

Untrusted sources, user turns, tools, external documents, and peer-agent messages enter the
\project{} adapter layer. The adapter converts each event into a canonical memory record. Write-time
admission either commits, down-trusts, or quarantines the record. The audit logger records every
decision. Retrieval calls fetch native candidates, then risk-aware reranking applies trust,
corroboration, contradiction, lineage, duplicate, and diversity logic. The context shaper renders
selected records as typed evidence blocks. The agent reasons over this shaped context. Before tools
run, pre-tool verification checks user authorization and trusted support. For multi-agent systems,
the propagation controller attaches lineage edges, decays trust across summaries, and freezes rapid
low-trust fan-out.
\end{minipage}}
\caption{Architecture placeholder. The final paper should replace or refine this with a clean
pipeline diagram showing trust boundaries and audit side channels.}
\end{figure}

\section{Components, Inputs, Outputs, and Dependencies}

\begin{longtable}{|L{0.18\textwidth}|L{0.22\textwidth}|L{0.22\textwidth}|L{0.28\textwidth}|}
\hline
\textbf{Component} & \textbf{Input} & \textbf{Output} & \textbf{Dependencies and notes} \\
\hline
Framework adapters & Native framework write, query, retrieve, add, and shared-memory calls. &
Canonical \texttt{MemoryRecord} and \texttt{MemoryEvent}. & One adapter each for LangGraph,
LlamaIndex, CrewAI. AutoGen optional. \\
\hline
Canonical schema & Raw content plus metadata. & Validated record with subsystem, source type,
actor, lineage, trust, labels, and timestamps. & Must remain identical across frameworks. \\
\hline
Provenance engine & Source type, actor, parent IDs, framework namespace, record labels. & Trust
prior and lineage features. & Uses configured trust priors and decay rules. \\
\hline
Write-time admission gate & Candidate write record. & Decision: commit, commit-low-trust,
quarantine, reject. & Uses instructionality, contradiction, sensitivity, duplicate, and source-risk
features. \\
\hline
Quarantine store & Suspicious record and decision rationale. & Non-retrievable or restricted
quarantine record. & Must be isolated from normal retrieval. \\
\hline
Audit logger & Every write, retrieve, rerank, context, tool, and propagation event. & Append-only
JSONL logs. & Required for reproducibility, metrics, forensics, and paper figures. \\
\hline
Retrieval reranker & Native top-k candidates, query, task context. & Risk-adjusted candidate set. &
Score combines relevance, trust, corroboration, contradiction, lineage risk, duplicate penalty, and
diversity. \\
\hline
Context shaper & Selected records and risk labels. & Typed evidence blocks. & Must prevent raw
retrieved text from being rendered as instruction. \\
\hline
Pre-tool verifier & Proposed tool call, user instruction, supporting memory records. & Allow,
ask-confirmation, read-only degrade, or block. & Uses support threshold and tool risk classes. \\
\hline
Propagation controller & Shared-memory writes, summaries, peer messages. & Lineage edges, trust
decay, freeze/quarantine decisions. & Needed for shared-memory cascade experiments. \\
\hline
Experiment runner & Scenario manifests, framework config, baseline config, seeds. & Run logs,
metrics, result tables, plots. & Must isolate runs and preserve raw logs. \\
\hline
\end{longtable}

\section{Canonical Record Schema}

The schema below is an implementation contract. The report gives the core fields; this manual adds
fields required for experiments, reproducibility, and lineage.

\begin{verbatim}
MemoryRecord:
  record_id: string
  run_id: string
  framework: enum[langgraph, llamaindex, crewai, autogen, letta]
  subsystem: enum[vector_rag, semantic, episodic, shared_coordination,
                  procedural_extension]
  content: string
  content_hash: string
  source_type: enum[system_seed, curated_corpus, user_text, external_doc,
                    tool_output, peer_agent, summary, synthetic_attack]
  source_id: string
  actor_id: string
  timestamp: iso8601
  parent_ids: list[string]
  lineage_depth: integer
  trust_prior: float
  trust_score: float
  relevance_score: float|null
  risk_score: float
  contradiction_score: float
  instructionality_score: float
  duplicate_cluster_id: string|null
  sensitivity_label: enum[none, private, credential_like, policy_like]
  poison_label: enum[benign, poisoned, unknown]
  decision: enum[commit, commit_low_trust, quarantine, reject]
  decision_reasons: list[string]
  quarantine_id: string|null
  tags: list[string]
\end{verbatim}

\missing{} The final implementation must decide whether records are persisted in SQLite, DuckDB,
JSONL plus vector index, or framework-native stores plus a sidecar metadata database. The quickest
MVP is JSONL audit logs plus SQLite metadata plus framework-native memory stores.

\section{Core API Contracts}

\begin{verbatim}
class MemoryFirewall:
    def admit_write(self, record: MemoryRecord) -> AdmissionDecision: ...
    def retrieve(self, query: MemoryQuery) -> RetrievalDecision: ...
    def shape_context(self, decision: RetrievalDecision) -> EvidenceBundle: ...
    def verify_tool_call(self, call: ToolCall,
                         evidence: EvidenceBundle) -> ToolDecision: ...
    def propagate(self, event: PropagationEvent) -> PropagationDecision: ...

class FrameworkAdapter:
    def wrap_memory(self, native_memory: object) -> object: ...
    def to_record(self, event: NativeMemoryEvent) -> MemoryRecord: ...
    def native_commit(self, record: MemoryRecord) -> None: ...
    def native_query(self, query: MemoryQuery) -> list[MemoryRecord]: ...
\end{verbatim}

The API should be synchronous first. Async support can be added behind the same interface only after
the synchronous MVP is stable.

\section{Scoring Logic}

The retrieval utility score should be explicit and logged:

\[
\begin{aligned}
S(r, q) ={}&
\lambda_{\mathrm{rel}}\mathrm{Rel}(r,q)
+ \lambda_{\mathrm{trust}}\mathrm{Trust}(r)
+ \lambda_{\mathrm{cor}}\mathrm{Corrob}(r)\\
&- \lambda_{\mathrm{contr}}\mathrm{Contradict}(r,q)
- \lambda_{\mathrm{lineage}}\mathrm{LineageRisk}(r)
- \lambda_{\mathrm{dup}}\mathrm{DuplicatePenalty}(r).
\end{aligned}
\]

The final implementation must log each term separately. If only the aggregate score is logged,
ablation and reviewer debugging become weak.

\missing{} Initial weights are not specified in the report. Start with fixed configuration values,
then run a threshold sensitivity sweep. Do not train a detector unless the MVP is already complete.

\section{Data Flow}

\begin{enumerate}
  \item A user turn, tool output, external document, or peer-agent message attempts to enter memory.
  \item The framework adapter converts the native event into a \texttt{MemoryRecord}.
  \item The provenance engine assigns a trust prior from source type and lineage depth.
  \item The write gate computes instructionality, contradiction, duplicate, sensitivity, and source
        risk features.
  \item The gate commits, down-trusts, quarantines, or rejects the record.
  \item The audit logger writes the raw event, normalized record, scores, and decision.
  \item During task execution, the native framework retrieves candidates.
  \item The retrieval gate reranks candidates and enforces source and lineage diversity.
  \item The context shaper renders memory as typed evidence blocks.
  \item The agent proposes an answer or tool call.
  \item The pre-tool verifier checks whether high-impact tool arguments are supported by explicit
        user intent or corroborated trusted memory.
  \item In multi-agent runs, shared-memory writes carry parent IDs and trust decay; rapid low-trust
        fan-out is frozen or quarantined.
\end{enumerate}

\section{Trust Boundaries and Attacker Interaction Points}

\begin{longtable}{|L{0.22\textwidth}|L{0.27\textwidth}|L{0.41\textwidth}|}
\hline
\textbf{Boundary} & \textbf{Attacker interaction point} & \textbf{Required control} \\
\hline
External content to memory & Poisoned documents, retrieved passages, tool-returned text. &
Source trust priors, instructionality scan, quarantine, context shaping. \\
\hline
User text to memory & Query-only write inducement and false fact insertion. & Write admission,
contradiction checks against high-trust memory, audit labels. \\
\hline
Memory to prompt context & Unsafe retrieved chunks entering the model as raw instructions. &
Risk-aware reranking and typed evidence rendering. \\
\hline
Memory to tool execution & Memory-conditioned tool argument or policy change. & Pre-tool
verification, support threshold, read-only degrade. \\
\hline
Agent to shared memory & Compromised peer-agent output, infected summaries, team notes. &
Lineage edges, trust decay, fan-out freeze. \\
\hline
Quarantine to normal memory & Accidental retrieval of suspicious records. & Physical namespace
separation and explicit quarantine policy. \\
\hline
\end{longtable}

\section{Implementation Order}

The first components to implement are the canonical schema, audit logger, LangGraph adapter,
write gate, retrieval gate, and metric calculation. Propagation should be implemented before any
optional fourth framework because propagation is one of the paper's differentiators.

\section{Required Paper Figures}

\begin{longtable}{|L{0.22\textwidth}|L{0.36\textwidth}|L{0.32\textwidth}|}
\hline
\textbf{Figure} & \textbf{What it should show} & \textbf{Why it matters} \\
\hline
System architecture & Pipeline from untrusted sources through gates to agent, tools, and shared
memory, with audit side channel and trust boundaries. & Establishes that the system is lifecycle
defense, not only retrieval filtering. \\
\hline
Canonical memory schema & Four memory subsystem examples normalized into one schema. & Supports
framework/subsystem portability. \\
\hline
Lineage graph & Poison source, summary hops, shared writes, and downstream actions. & Makes
propagation measurable and concrete. \\
\hline
Security-cost frontier & Attack reduction or contaminated context reduction versus p95 latency or
memory overhead. & Preempts ``too expensive'' reviewer objections. \\
\hline
Scalability curves & Memory size, poison ratio, top-k, and agent count sweeps. & Makes scaling a
central contribution rather than appendix filler. \\
\hline
\end{longtable}

\chapter{Threat Model}

\section{Attacker Capabilities}

The attacker has bounded but realistic capabilities:

\begin{itemize}
  \item \textbf{Query-only write inducement:} the attacker can interact with the agent through normal
        user or task channels and attempt to cause persistent malicious or misleading memory writes.
  \item \textbf{Corpus or passage poisoning:} the attacker can insert or influence documents likely
        to be retrieved by vector/RAG memory.
  \item \textbf{Indirect prompt injection:} the attacker can place instructions inside retrieved or
        tool-returned content.
  \item \textbf{Compromised peer-agent output:} the attacker can control or contaminate one peer
        agent's output that is written into shared coordination memory.
\end{itemize}

The attacker does not control the model weights, host operating system, memory firewall code, or
experiment harness.

\section{Attacker Goals}

\begin{itemize}
  \item Cause poisoned records to be committed to durable memory.
  \item Cause poisoned or low-trust records to dominate retrieval.
  \item Cause retrieved content to enter reasoning as instruction rather than evidence.
  \item Cause contradiction collisions where recent or semantically similar false records override
        high-trust records.
  \item Cause memory-conditioned unsafe tool calls.
  \item Cause contamination to propagate across agents through shared memory, summaries, or notes.
\end{itemize}

\section{Defender Assumptions}

\begin{itemize}
  \item The defender can wrap or intercept framework memory APIs.
  \item The defender can maintain sidecar metadata, audit logs, and quarantine namespaces.
  \item The defender can assign source trust priors based on source type and framework context.
  \item The defender can run lightweight deterministic checks and optional LLM checks.
  \item The defender can enforce tool-call confirmation or read-only degradation for high-risk calls.
\end{itemize}

\missing{} The implementation must decide whether optional LLM checks use the same model as the
agent, a cheaper verifier model, or no model at all.

\section{In Scope}

\begin{itemize}
  \item Persistent memory contamination through normal interaction paths.
  \item Vector/RAG poisoning that affects retrieval.
  \item Semantic false fact insertion and contradiction.
  \item Episodic poisoned exemplar replay.
  \item Shared coordination memory cascades.
  \item Memory flooding and duplicate domination.
  \item Utility and latency impact of defensive gates.
\end{itemize}

\section{Out of Scope}

\begin{itemize}
  \item Model-weight poisoning or fine-tuning attacks.
  \item Host compromise, file-system compromise, or arbitrary code execution.
  \item Embedding-model attacks and vector-index implementation vulnerabilities.
  \item Vendor prompt leakage or hidden system prompt extraction.
  \item Covert channels.
  \item Human-in-the-loop social engineering outside the agent interface.
  \item Formal proof of complete memory integrity.
\end{itemize}

\section{Victim System and Environment}

The victim is an LLM-agent application built on a public framework with durable memory enabled. The
MVP victim systems are:

\begin{itemize}
  \item LangGraph agent with vector/RAG and semantic memory.
  \item LlamaIndex agent with memory blocks or long-term memory abstractions.
  \item CrewAI team with shared coordination memory.
\end{itemize}

All tasks should use safe stub tools and synthetic or public non-sensitive corpora. The paper should
avoid regulated domains unless there is a clear reason and proper ethics review.

\section{Credibility of the Threat Model}

The threat model is credible because it matches the direction of public agent systems: persistent
memory, retrieval, summarization, and multi-agent coordination are exposed as normal features. It is
not exaggerated because it excludes host compromise, model compromise, hidden vendor internals, and
unbounded attacker control. The attacker interacts through plausible content and memory paths.

\chapter{Implementation Blueprint}

\section{Repository Structure}

Create the following repository structure:

\begin{verbatim}
memsieve/
  pyproject.toml
  README.md
  LICENSE
  configs/
    default.yaml
    frameworks.yaml
    thresholds.yaml
    models.yaml
  memsieve/
    __init__.py
    core/
      record.py
      events.py
      decisions.py
      scoring.py
      config.py
      clock.py
    stores/
      metadata_store.py
      quarantine_store.py
      audit_log.py
      vector_store.py
    gates/
      write_gate.py
      retrieval_gate.py
      context_shaper.py
      tool_verifier.py
      propagation_controller.py
    adapters/
      base.py
      langgraph_adapter.py
      llamaindex_adapter.py
      crewai_adapter.py
      autogen_adapter.py
    baselines/
      no_defense.py
      provenance_only.py
      spotlighting_wrapper.py
      trust_threshold.py
      random_quarantine.py
    scenarios/
      poisoned_fact.py
      retrieved_instruction.py
      experience_graft.py
      conflict_recency.py
      shared_cascade.py
      flooding.py
    experiments/
      runner.py
      manifests.py
      seed.py
      graders.py
      metrics.py
      aggregate.py
    plotting/
      tables.py
      figures.py
      latex_export.py
  manifests/
    scenarios/
    experiment_matrix.yaml
  data/
    raw/
    synthetic/
    curated/
  runs/
    .gitkeep
  scripts/
    run_smoke_tests.ps1
    run_main_matrix.ps1
    aggregate_results.ps1
    make_figures.ps1
  tests/
    unit/
    integration/
    regression/
  paper/
    memsieve_paper.tex
    figures/
    tables/
\end{verbatim}

\section{Core Modules and Classes}

\begin{longtable}{|L{0.25\textwidth}|L{0.30\textwidth}|L{0.35\textwidth}|}
\hline
\textbf{Module} & \textbf{Core classes} & \textbf{Responsibility} \\
\hline
\texttt{core.record} & \texttt{MemoryRecord}, \texttt{MemoryQuery} & Canonical schema,
validation, serialization. \\
\hline
\texttt{core.decisions} & \texttt{AdmissionDecision}, \texttt{RetrievalDecision},
\texttt{ToolDecision}, \texttt{PropagationDecision} & Typed decision objects with scores and
reasons. \\
\hline
\texttt{stores.audit\_log} & \texttt{AuditLogger} & Append-only JSONL event logging. \\
\hline
\texttt{stores.quarantine\_store} & \texttt{QuarantineStore} & Isolated suspicious record store. \\
\hline
\texttt{gates.write\_gate} & \texttt{WriteGate} & Write scoring, contradiction checks, duplicate
throttling, quarantine decisions. \\
\hline
\texttt{gates.retrieval\_gate} & \texttt{RetrievalGate} & Risk-aware reranking and diversity
constraints. \\
\hline
\texttt{gates.context\_shaper} & \texttt{ContextShaper} & Converts raw records into typed evidence
blocks. \\
\hline
\texttt{gates.tool\_verifier} & \texttt{ToolVerifier} & Blocks, confirms, or degrades risky
memory-conditioned tool calls. \\
\hline
\texttt{gates.propagation\_controller} & \texttt{PropagationController} & Adds lineage, decays
trust, detects fan-out cascades. \\
\hline
\texttt{adapters.*} & \texttt{FrameworkAdapter} subclasses & Wrap native memory APIs. \\
\hline
\texttt{experiments.runner} & \texttt{ExperimentRunner} & Runs scenario matrix and writes
structured artifacts. \\
\hline
\texttt{experiments.metrics} & \texttt{MetricComputer} & Computes ASR, URR, CCR, BUP, TSD,
propagation, latency, and scaling metrics. \\
\hline
\end{longtable}

\section{Configuration Files}

\begin{verbatim}
# configs/default.yaml
run:
  seed: 0
  output_dir: runs
  log_level: INFO

models:
  agent_model: MISSING_MUST_DECIDE
  verifier_model: none
  embedding_model: MISSING_MUST_DECIDE

thresholds:
  quarantine_risk: 0.75
  low_trust_risk: 0.45
  contradiction_high: 0.70
  instructionality_high: 0.65
  propagation_freeze_fanout: 3
  trust_decay_per_hop: 0.15

retrieval:
  top_k_native: 16
  top_k_final: 4
  require_source_diversity: true
  max_same_lineage: 1

tool_verification:
  high_impact_tools: [write_file, send_message, external_post, purchase]
  require_user_or_trusted_support: true
\end{verbatim}

The threshold values above are starting defaults, not scientific facts. The evaluation must include a
sensitivity sweep.

\section{Logging Design}

Every run must write:

\begin{itemize}
  \item \texttt{run\_manifest.yaml}: framework, model, seed, scenario, baseline, thresholds, git hash
        if available, and environment.
  \item \texttt{events.jsonl}: every normalized memory event.
  \item \texttt{records.jsonl}: canonical records and labels.
  \item \texttt{decisions.jsonl}: gate scores and decisions.
  \item \texttt{retrievals.jsonl}: native candidates, reranked candidates, selected context.
  \item \texttt{tool\_calls.jsonl}: proposed and final tool decisions.
  \item \texttt{lineage.jsonl}: parent-child edges for propagation.
  \item \texttt{grader.jsonl}: task success, attack success, and grading rationale.
  \item \texttt{timing.jsonl}: p50, p95, and per-stage latencies.
\end{itemize}

\textbf{Hard rule:} paper metrics must be computed from logs, not from ad hoc notebook state.

\section{Experiment Runner Design}

The runner should execute a matrix:

\begin{verbatim}
for framework in frameworks:
  for subsystem in subsystems:
    for scenario in scenarios:
      for defense in defenses:
        for seed in seeds:
          prepare_clean_environment()
          seed_memory()
          inject_attack_if_needed()
          run_task()
          grade_outputs()
          save_logs()
\end{verbatim}

Each run directory must be immutable after completion. Re-running should create a new run ID.

\section{Reproducibility Design}

\begin{itemize}
  \item Pin Python version and all framework versions in \texttt{pyproject.toml} or a lock file.
  \item Save all scenario manifests and generated synthetic records.
  \item Save model names, API parameters, temperature, top-p, max tokens, and embedding model.
  \item Use deterministic stub tools for all side-effectful actions.
  \item Use fixed seeds for synthetic generation, retrieval tie-breaking, and scenario ordering.
  \item Export paper tables directly from result JSON or CSV.
  \item Keep raw logs for every table row.
\end{itemize}

\missing{} Exact reproducibility policy for hosted models: the paper must disclose that hosted LLMs
can change unless a pinned snapshot or local model is used.

\section{Implementation Phases}

\subsection{Phase 0: Scaffold and Reproducibility}

\textbf{Goal.} Create the repository, configuration, logging, and test skeleton.

\textbf{Required modules.} \texttt{core.config}, \texttt{stores.audit\_log},
\texttt{experiments.runner}.

\textbf{Exact tasks.}
\begin{itemize}
  \item Create package structure and dependency file.
  \item Implement config loading and schema validation.
  \item Implement immutable run directory creation.
  \item Implement append-only JSONL writer.
  \item Add unit tests for config and logging.
\end{itemize}

\textbf{Expected outputs.} A smoke run creates a run directory with manifest and empty logs.

\textbf{Correctness checks.} Tests pass; invalid config fails loudly; logs are valid JSONL.

\textbf{Common failure points.} Silent overwrites, unpinned configs, invalid JSON due to raw model
text.

\textbf{Evidence to save.} Smoke-test logs and configuration snapshot.

\subsection{Phase 1: Canonical Schema and Core Decisions}

\textbf{Goal.} Implement canonical records and typed decisions.

\textbf{Required modules.} \texttt{core.record}, \texttt{core.events},
\texttt{core.decisions}.

\textbf{Exact tasks.}
\begin{itemize}
  \item Implement \texttt{MemoryRecord} validation.
  \item Implement content hashing and deterministic record IDs.
  \item Implement admission, retrieval, tool, and propagation decision classes.
  \item Add serialization round-trip tests.
\end{itemize}

\textbf{Expected outputs.} Any framework event can be represented as a stable record.

\textbf{Correctness checks.} Round-trip tests preserve all fields; missing required fields fail.

\textbf{Common failure points.} Framework-specific fields leak into generic schema; IDs are
non-deterministic.

\textbf{Evidence to save.} Schema examples for all four subsystems.

\subsection{Phase 2: Framework Adapter Smoke Tests}

\textbf{Goal.} Wrap LangGraph, LlamaIndex, and CrewAI enough to intercept memory writes and
retrievals.

\textbf{Required modules.} \texttt{adapters.base}, \texttt{adapters.langgraph\_adapter},
\texttt{adapters.llamaindex\_adapter}, \texttt{adapters.crewai\_adapter}.

\textbf{Exact tasks.}
\begin{itemize}
  \item Identify native memory write and retrieval calls.
  \item Implement adapter methods for write, query, and commit.
  \item Build one clean task per framework.
  \item Verify identical event logs across framework-specific runs.
\end{itemize}

\textbf{Expected outputs.} Three clean smoke runs with memory event logs.

\textbf{Correctness checks.} Each adapter captures writes and retrievals without changing clean
answers.

\textbf{Common failure points.} Framework version drift, async behavior, hidden memory writes, and
agent code bypassing the adapter.

\textbf{Evidence to save.} Adapter smoke logs and minimal working examples.

\subsection{Phase 3: Write Gate, Quarantine, and Audit}

\textbf{Goal.} Stop or label contaminated writes before they enter normal memory.

\textbf{Required modules.} \texttt{gates.write\_gate}, \texttt{stores.quarantine\_store},
\texttt{stores.metadata\_store}.

\textbf{Exact tasks.}
\begin{itemize}
  \item Implement trust priors by source type.
  \item Implement heuristic instructionality features.
  \item Implement contradiction checks against high-trust semantic memory.
  \item Implement near-duplicate clustering or hashing.
  \item Implement quarantine namespace and audit reasons.
\end{itemize}

\textbf{Expected outputs.} Poisoned writes are committed, down-trusted, quarantined, or rejected with
logged reasons.

\textbf{Correctness checks.} Unit tests for benign writes, obvious poisoned writes, contradictions,
and duplicate floods.

\textbf{Common failure points.} Over-blocking benign user preferences; quarantined records leaking
into retrieval.

\textbf{Evidence to save.} Quarantine precision/recall logs and example decision rationales.

\subsection{Phase 4: Retrieval Gate and Context Shaping}

\textbf{Goal.} Make retrieval optimize task utility under risk, not raw similarity alone.

\textbf{Required modules.} \texttt{gates.retrieval\_gate}, \texttt{gates.context\_shaper}.

\textbf{Exact tasks.}
\begin{itemize}
  \item Retrieve a larger native candidate pool.
  \item Compute risk-aware score terms.
  \item Enforce source and lineage diversity.
  \item Render records as trusted facts, examples, shared notes, contradictions, or quarantined
        excerpts.
  \item Log native rank, defended rank, selected context, and dropped candidates.
\end{itemize}

\textbf{Expected outputs.} The model sees typed evidence blocks and fewer contaminated records.

\textbf{Correctness checks.} Retrieval tests where poisoned top-1 native result is demoted;
context output contains no untyped raw injected instruction.

\textbf{Common failure points.} Excessive relevance loss; context shaper accidentally quotes
malicious content as command-like text.

\textbf{Evidence to save.} Unsafe Retrieval Rate and Contaminated Context Rate examples.

\subsection{Phase 5: Pre-Tool Verification and Propagation Control}

\textbf{Goal.} Prevent memory-conditioned unsafe actions and shared-memory cascades.

\textbf{Required modules.} \texttt{gates.tool\_verifier}, \texttt{gates.propagation\_controller}.

\textbf{Exact tasks.}
\begin{itemize}
  \item Define high-impact stub tools.
  \item Require explicit user instruction or trusted corroborating memory for high-impact actions.
  \item Attach parent IDs to shared-memory writes.
  \item Decay trust across summaries and peer-agent hops.
  \item Freeze rapid low-trust fan-out.
\end{itemize}

\textbf{Expected outputs.} Tool calls and shared writes have support evidence and lineage.

\textbf{Correctness checks.} A compromised peer cannot contaminate all downstream agents in the
shared-memory cascade scenario.

\textbf{Common failure points.} Tool verifier blocks too many legitimate actions; lineage is missing
for summaries.

\textbf{Evidence to save.} Lineage graph, propagation depth, fan-out, and blocked-call logs.

\subsection{Phase 6: Baselines, Scenarios, and Graders}

\textbf{Goal.} Implement fair comparisons and auto-gradable scenarios.

\textbf{Required modules.} \texttt{baselines.*}, \texttt{scenarios.*}, \texttt{experiments.graders}.

\textbf{Exact tasks.}
\begin{itemize}
  \item Implement no defense, provenance-only, spotlighting-style wrapper, trust-threshold sanitizer,
        and random quarantine baseline.
  \item Implement six scenario families.
  \item Implement deterministic graders for task success and attack success.
  \item Add scenario manifests with seed, memory subsystem, attack records, and expected labels.
\end{itemize}

\textbf{Expected outputs.} A small matrix can run end to end and produce metrics.

\textbf{Correctness checks.} Baselines use the same tasks, models, memory sizes, and seeds as
\project{}.

\textbf{Common failure points.} Baselines accidentally get weaker retrieval settings; graders rely
on brittle string matching only.

\textbf{Evidence to save.} Scenario manifests and grader validation examples.

\subsection{Phase 7: Full Evaluation Pipeline}

\textbf{Goal.} Run the MVP experiment matrix and produce paper-ready CSV, LaTeX tables, and plots.

\textbf{Required modules.} \texttt{experiments.metrics}, \texttt{experiments.aggregate},
\texttt{plotting.tables}, \texttt{plotting.figures}.

\textbf{Exact tasks.}
\begin{itemize}
  \item Run main attack matrix.
  \item Run clean utility matrix.
  \item Run ablations.
  \item Run memory-size, poison-ratio, top-k, and agent-count sweeps.
  \item Aggregate metrics with confidence intervals.
  \item Export plots and tables.
\end{itemize}

\textbf{Expected outputs.} Paper tables and figures are generated from raw logs.

\textbf{Correctness checks.} Re-running aggregation from raw logs reproduces table values.

\textbf{Common failure points.} Missing run metadata, inconsistent seeds, incomplete failed-run
handling.

\textbf{Evidence to save.} Raw logs, aggregate CSVs, plotting scripts, figure source files.

\chapter{Scenario Design}

\section{Scenario Requirements}

Every scenario must define: name, purpose, initial system state, attacker action, defender behavior,
expected failure or success mode, metrics, logs, table/figure output, and supported claim.

\section{Scenario Matrix}

\begin{longtable}{|L{0.14\textwidth}|L{0.18\textwidth}|L{0.18\textwidth}|L{0.20\textwidth}|L{0.20\textwidth}|}
\hline
\textbf{Scenario} & \textbf{Purpose} & \textbf{Attacker action} & \textbf{Defender behavior} &
\textbf{Claim supported} \\
\hline
Poisoned fact insertion & Test persistent false semantic memory. & Induce storage of a false
project or user fact through low-trust content. & Quarantine, down-trust, or mark contradiction
against high-trust fact. & Write-path defense reduces poisoned persistent memory. \\
\hline
Retrieved instruction hijack & Test RAG chunks that contain hidden instructions. & Insert a
retrievable document that tells the agent to change behavior. & Demote, quarantine, or render as
untrusted evidence. & Retrieval and context shaping reduce unsafe retrieved instruction. \\
\hline
Experience graft & Test poisoned episodic examples. & Seed an apparently successful but unsafe
trajectory. & Suppress or label the exemplar; prefer trusted examples. & Episodic memory needs
defense, not only RAG. \\
\hline
Conflict and recency collision & Test recent false record versus older trusted record. & Add recent
semantically overlapping contradiction. & Surface contradiction and prefer high-trust record or
abstain. & Trust and contradiction beat naive recency/similarity. \\
\hline
Shared-memory cascade & Test peer-agent contamination. & Compromised peer writes contaminated
summary or shared note. & Attach lineage, decay trust, freeze low-trust fan-out. & Propagation
controller caps cross-agent spread. \\
\hline
Memory flooding & Test duplicate and top-k domination. & Insert many near-duplicate low-quality or
poisoned records. & Cluster duplicates, throttle writes, enforce diversity. & Defense remains stable
under memory pressure. \\
\hline
\end{longtable}

\section{Detailed Test Cases}

\subsection{Poisoned Fact Insertion}

\textbf{System state before experiment.} Semantic memory contains high-trust seeded facts and benign
project/user facts. The agent has a task requiring later recall of one fact.

\textbf{Attacker action.} The attacker submits low-trust content that implies a false persistent fact
should be saved.

\textbf{Defender behavior.} The write gate computes source risk, contradiction score, and
instructionality. It commits benign facts, down-trusts uncertain facts, and quarantines high-risk
false records.

\textbf{Expected failure mode without defense.} The false fact is written, retrieved later, and used
in an answer or plan.

\textbf{Metrics.} Attack Success Rate, Quarantine Precision/Recall, Contaminated Context Rate, Task
Success Under Defense.

\textbf{Logs needed.} Write event, scores, decision reasons, retrieval event, final answer, grader.

\textbf{Output.} Main attack table and write-gate ablation row.

\textbf{Claim supported.} \project{} reduces persistent poisoned writes while preserving clean
semantic memory utility.

\subsection{Retrieved Instruction Hijack}

\textbf{System state before experiment.} Vector/RAG memory contains benign documents relevant to a
question. Attack document is semantically relevant and likely to appear in native top-k.

\textbf{Attacker action.} The attacker inserts a poisoned document with hidden or explicit
instructions embedded in content.

\textbf{Defender behavior.} Retrieval gate scores lineage and instructionality risk, applies source
diversity, and context shaper renders the chunk as untrusted evidence rather than instruction.

\textbf{Expected failure mode without defense.} The model follows the retrieved instruction or
changes the answer/tool plan.

\textbf{Metrics.} Unsafe Retrieval Rate, Contaminated Context Rate, Attack Success Reduction, BUP,
latency overhead.

\textbf{Logs needed.} Native top-k, defended top-k, context bundle, final output, stage timings.

\textbf{Output.} Unsafe retrieval table and memory-size scaling curve.

\textbf{Claim supported.} The system improves over raw vector retrieval and spotlighting-style
context wrapping.

\subsection{Experience Graft}

\textbf{System state before experiment.} Episodic memory contains prior successful trajectories. A
task requires selecting or imitating a prior example.

\textbf{Attacker action.} The attacker seeds a poisoned trajectory marked as successful but containing
an unsafe step.

\textbf{Defender behavior.} The write gate marks low trust or quarantines the trajectory; retrieval
gate penalizes unsupported low-trust exemplars; context shaper labels examples as examples, not
rules.

\textbf{Expected failure mode without defense.} The agent imitates the poisoned plan.

\textbf{Metrics.} Attack Success Rate, selected exemplar poison rate, TSD, BUP.

\textbf{Logs needed.} Episodic write labels, selected exemplars, generated plan, tool sequence.

\textbf{Output.} Non-RAG subsystem row in main table and ablation table.

\textbf{Claim supported.} Memory defense generalizes beyond vector/RAG corpora.

\subsection{Conflict and Recency Collision}

\textbf{System state before experiment.} Semantic or shared memory contains an older high-trust
record. A later low-trust contradictory record overlaps semantically.

\textbf{Attacker action.} Insert a recent contradiction likely to win naive recency or similarity
ranking.

\textbf{Defender behavior.} Contradiction detector flags conflict; retrieval gate prefers high-trust
record or surfaces both with a warning; tool verifier blocks actions relying only on the low-trust
branch.

\textbf{Expected failure mode without defense.} The recent false record silently overrides trusted
memory.

\textbf{Metrics.} Contradiction win rate, TSD, CCR, abstention/confirmation rate.

\textbf{Logs needed.} Contradiction pairs, ranking, context, final answer/action.

\textbf{Output.} Conflict-resolution table.

\textbf{Claim supported.} Provenance and contradiction signals improve over similarity/recency.

\subsection{Shared-Memory Cascade}

\textbf{System state before experiment.} CrewAI team has shared coordination memory and multiple
agents. One peer is contaminated.

\textbf{Attacker action.} The compromised peer writes a contaminated summary or shared note.

\textbf{Defender behavior.} Propagation controller records lineage, decays trust, freezes rapid
low-trust fan-out, and quarantines suspicious shared writes.

\textbf{Expected failure mode without defense.} Contamination spreads to downstream agents and
affects their writes or tool actions.

\textbf{Metrics.} Propagation depth, fan-out, contaminated shared writes, TSD, latency overhead.

\textbf{Logs needed.} Lineage graph, shared writes, downstream retrievals, tool calls.

\textbf{Output.} Agent-count scaling figure and lineage graph.

\textbf{Claim supported.} \project{} contains memory-mediated cross-agent spread.

\subsection{Memory Flooding}

\textbf{System state before experiment.} Memory store contains benign records. Attack condition
inserts many near-duplicate records with low trust or poison labels.

\textbf{Attacker action.} Flood vector, episodic, or shared memory with semantically similar records.

\textbf{Defender behavior.} Duplicate throttling, source diversity, and max-same-lineage constraints
prevent domination.

\textbf{Expected failure mode without defense.} Top-k dominated by attacker cluster; utility drops;
latency rises.

\textbf{Metrics.} Top-k Source Diversity, URR, CCR, BUP, latency overhead, memory growth overhead.

\textbf{Logs needed.} Duplicate cluster IDs, retrieval source distribution, timings, task outcomes.

\textbf{Output.} Flooding stress table and security-cost frontier.

\textbf{Claim supported.} Defense remains stable under adversarial memory pressure.

\chapter{Experimental Methodology}

\section{Experimental Setup}

The MVP setup uses three public frameworks, four memory subsystems, six scenarios, five defense
settings, three seeds, one primary model, and one shared embedding backend.

\begin{itemize}
  \item Frameworks: LangGraph, LlamaIndex, CrewAI.
  \item Optional: AutoGen after MVP results are complete.
  \item Memory subsystems: vector/RAG, semantic, episodic, shared coordination.
  \item Defense settings: no defense, provenance-only, spotlighting-style retrieval wrapper,
        trust-threshold sanitizer, full \project{}.
  \item Seeds: at least 3 for scenario generation and retrieval tie-breaking.
  \item Model temperature: low or zero for reproducibility.
\end{itemize}

\missing{} Exact hosted/local model, model version, embedding backend, and API parameters.

\section{Datasets and Simulated Environment}

Use synthetic and public non-sensitive corpora. The safest MVP is a synthetic project-assistant
environment with:

\begin{itemize}
  \item project documents for vector/RAG memory;
  \item user and project facts for semantic memory;
  \item prior task trajectories for episodic memory;
  \item team notes and summaries for shared coordination memory;
  \item deterministic stub tools for high-impact actions.
\end{itemize}

This avoids domain-specific ethics drag while still exercising real memory paths.

\missing{} The report does not specify a public corpus. Choose one public, permissively licensed
corpus only if it adds realism without slowing the MVP.

\section{Baselines}

Run every applicable scenario under:

\begin{itemize}
  \item No defense.
  \item Provenance-only filter.
  \item Spotlighting-style context wrapper.
  \item Trust-threshold sanitizer.
  \item Random quarantine baseline.
  \item Full \project{}.
\end{itemize}

Optional comparisons:

\begin{itemize}
  \item RAGForensics for vector/RAG rows if code is easy to reuse.
  \item Task Shield for tool-action rows if integration is moderate.
\end{itemize}

\section{Ablations}

\begin{itemize}
  \item Full \project{} minus write gate.
  \item Full \project{} minus retrieval risk scoring.
  \item Full \project{} minus context shaping.
  \item Full \project{} minus contradiction checks.
  \item Full \project{} minus source diversity.
  \item Full \project{} minus pre-tool verification.
  \item Full \project{} minus propagation controller.
  \item Full \project{} with no trust decay.
\end{itemize}

\section{Main Experiment Matrix}

\begin{longtable}{|L{0.17\textwidth}|L{0.22\textwidth}|L{0.19\textwidth}|L{0.22\textwidth}|L{0.10\textwidth}|}
\hline
\textbf{Experiment} & \textbf{Frameworks} & \textbf{Subsystems} & \textbf{Primary metrics} &
\textbf{Priority} \\
\hline
Clean utility & LangGraph, LlamaIndex, CrewAI & All four & BUP, TSD, latency & Must \\
\hline
Write-path poisoning & LangGraph, LlamaIndex & Semantic, episodic & ASR reduction, quarantine P/R
& Must \\
\hline
Unsafe retrieval & All MVP frameworks & Vector/RAG, semantic & URR, CCR, TSD & Must \\
\hline
Conflict resolution & LangGraph, CrewAI & Semantic, shared & Contradiction win rate, TSD & Must \\
\hline
Shared cascade & CrewAI, optional AutoGen & Shared coordination & Propagation depth, fan-out, TSD
& Must \\
\hline
Flooding stress & All MVP frameworks & Vector, episodic, shared & Top-k diversity, URR, latency &
Must \\
\hline
Memory-size scaling & All MVP frameworks & Applicable subsystems & Scalability slope, URR, CCR,
latency & Must \\
\hline
Agent-count scaling & CrewAI, optional AutoGen & Shared coordination & Propagation depth, fan-out,
latency & Must \\
\hline
Second-model transfer & One or two frameworks & Vector plus shared & ASR reduction, BUP & Optional
\\
\hline
\end{longtable}

\section{Repeated Trials and Seeds}

Use at least three seeds. For stochastic hosted models, run more if budget allows. Report confidence
intervals for aggregate rates. If only three seeds are feasible, do not overstate statistical
significance.

\section{Clean, Attacked, and Defended Conditions}

\begin{itemize}
  \item \textbf{Clean:} no poisoned records; measure utility, latency, memory overhead.
  \item \textbf{Attacked/no-defense:} poisoned records or compromised peer present; no gates.
  \item \textbf{Attacked/baseline:} same attack under each baseline.
  \item \textbf{Attacked/full defense:} same attack under full \project{}.
  \item \textbf{Attacked/ablation:} same attack with one module removed.
\end{itemize}

All comparisons must use identical tasks, seeds, framework versions, memory sizes, top-k values, and
model parameters.

\section{Tables and Figures to Produce}

\begin{itemize}
  \item Table 1: Frameworks, memory subsystems, adapter coverage.
  \item Table 2: Scenario matrix and threat capability mapping.
  \item Table 3: Main attack reduction results.
  \item Table 4: Clean utility preservation.
  \item Table 5: Ablation results.
  \item Table 6: Baseline comparison by subsystem.
  \item Figure 1: System architecture.
  \item Figure 2: Unsafe retrieval and contaminated context versus memory size.
  \item Figure 3: Cross-agent propagation depth versus agent count.
  \item Figure 4: Security-cost frontier.
  \item Figure 5: Threshold sensitivity or top-k source diversity.
\end{itemize}

\chapter{Metrics and Evaluation Logic}

\section{Notation}

Let $Q_a$ be attack queries, $Q_b$ be benign queries, $P$ be poisoned records, $R(q)$ be the final
retrieved set for query $q$, and $C(q)$ be the final memory-conditioned context shown to the model.
Let $U$ be task utility and $T$ be end-to-end latency.

\section{Metric Definitions}

\begin{longtable}{|L{0.16\textwidth}|L{0.24\textwidth}|L{0.20\textwidth}|L{0.20\textwidth}|L{0.10\textwidth}|}
\hline
\textbf{Metric} & \textbf{Formal definition} & \textbf{Intuition and computation} &
\textbf{Supports / weakens claim} & \textbf{Range} \\
\hline
Attack Success Rate (ASR) & $\frac{1}{|Q_a|}\sum_{q\in Q_a}\mathbb{1}[\mathrm{attack\ succeeds}]$
& Fraction of attacked tasks where the attacker achieves the scenario goal. Compute from grader
logs. & Lower ASR supports security; unchanged ASR weakens the paper. & $[0,1]$ \\
\hline
Attack Success Reduction & $1-\frac{\mathrm{ASR}_{def}}{\mathrm{ASR}_{base}}$ & Relative reduction
against no-defense or baseline. & Large positive reduction supports headline claim; negative values
mean defense is harmful. & $(-\infty,1]$ \\
\hline
Unsafe Retrieval Rate (URR) & $\frac{1}{|Q_a|}\sum_{q\in Q_a}\mathbb{1}[R(q)\cap P\neq\emptyset]$
& Measures whether poisoned records remain in final retrieval. Requires retrieved IDs and poison
labels. & Supports retrieval safety but is insufficient alone. & $[0,1]$ \\
\hline
Contaminated Context Rate (CCR) &
$\frac{1}{|Q|}\sum_q\mathbb{1}[\mathrm{contaminated}(C(q))]$ &
Measures what the model actually sees after context shaping. & Stronger than URR; high CCR weakens
context-shaping claims. & $[0,1]$ \\
\hline
Propagation Depth & Maximum or mean graph distance from initial poisoned write to downstream
contaminated writes/actions. & Build from parent-child lineage edges. & Lower depth supports
containment; depth growing with agent count weakens claim. & $\geq 0$ \\
\hline
Propagation Fan-out & Count of downstream contaminated records/actions reachable from initial
poison. & Compute from lineage DAG. & Lower fan-out supports shared-memory defense. & $\geq 0$ \\
\hline
Quarantine Precision & $\frac{TP}{TP+FP}$ & Fraction of quarantined records that were truly poisoned.
Requires poison labels. & Low precision means over-blocking. & $[0,1]$ \\
\hline
Quarantine Recall & $\frac{TP}{TP+FN}$ & Fraction of poisoned writes caught by quarantine or
down-trust policy. & Low recall means poisoned writes pass through. & $[0,1]$ \\
\hline
Benign Utility Preservation (BUP) & $\frac{U^{clean}_{def}}{U^{clean}_{base}}$ & Clean-task success
ratio under defense versus baseline. & Values near 1 support deployability; large drops are fatal. &
$[0,\infty)$ \\
\hline
Task Success Under Defense (TSD) & $U^{attack}_{def}$ & Utility on attacked tasks while defense is
enabled. & Security is less meaningful if the agent stops solving tasks. & Task-specific \\
\hline
Latency Overhead & $\frac{T_{def}-T_{base}}{T_{base}}$ at p50 and p95 & End-to-end and per-stage
timing. & Moderate overhead supports practicality; high p95 weakens systems claim. & $\geq -1$ \\
\hline
Memory Growth Overhead & $\frac{S_{def}-S_{base}}{S_{base}}$ & Extra records/bytes from metadata,
logs, quarantine, and lineage. & Supports storage scalability if bounded. & $\geq -1$ \\
\hline
Scalability Slope & $\hat{\beta}$ in $Y=\alpha+\beta\log N$ for $Y\in\{\mathrm{latency}, URR, CCR\}$
& Fit slopes over memory records, agents, or top-k. & Flatter security degradation supports scaling.
& Real-valued \\
\hline
Top-k Source Diversity & $\frac{|\mathrm{unique\ sources\ in\ }R(q)|}{|R(q)|}$ averaged over queries
& Measures whether retrieval is dominated by one source or lineage. & Higher diversity supports
flooding resistance. & $[0,1]$ \\
\hline
Contradiction Handling Rate & Fraction of conflict trials where the trusted branch is selected,
both branches are surfaced, or the agent abstains. & Requires labeled contradiction pairs. & High
rate supports semantic-memory safety. & $[0,1]$ \\
\hline
\end{longtable}

\section{Are New Metrics Needed?}

Yes. Existing attack success and utility metrics are necessary but not enough. Unsafe Retrieval Rate
is common-sense and easy to compute, but it misses context shaping: a poisoned record can be
retrieved and safely rendered as untrusted evidence. Contaminated Context Rate is therefore needed.
Propagation Depth and Fan-out are also needed because the report's distinctive claim involves
shared-memory containment. Scalability Slope is needed because the paper's thesis says security must
survive memory growth and team growth.

\chapter{Baselines and Ablations}

\section{Baseline Definitions}

\begin{longtable}{|L{0.20\textwidth}|L{0.34\textwidth}|L{0.26\textwidth}|}
\hline
\textbf{Baseline} & \textbf{Definition} & \textbf{Why fair} \\
\hline
No defense & Native framework memory with no additional filtering, reranking, context shaping, or
propagation control. & Necessary lower bound and deployment default. \\
\hline
Provenance-only & Attach source labels and trust priors but do not quarantine, rerank by risk, or
block tools. & Tests whether metadata alone explains results. \\
\hline
Spotlighting-style wrapper & Separate retrieved content from instructions and mark provenance in
prompt context. & Fair retrieval-context baseline named by the report. \\
\hline
Trust-threshold sanitizer & Drop or sanitize records below a trust threshold. & Strong lightweight
baseline close to practical defenses. \\
\hline
Random quarantine & Quarantine the same number of records as \project{} at random. & Detects whether
improvement comes merely from blocking volume. \\
\hline
RAGForensics optional & Vector/RAG-specific traceback defense if easy to integrate. & Fair only for
vector/RAG rows. Do not use for semantic/shared claims. \\
\hline
Task Shield optional & Tool-time task-alignment defense if easy to integrate. & Fair only for
tool-action rows. Do not use for write-path claims. \\
\hline
\end{longtable}

\section{Comparative Evaluation Logic}

A convincing comparison is not ``\project{} beats no defense.'' The paper must show:

\begin{itemize}
  \item \project{} beats no defense on ASR, URR, CCR, and propagation.
  \item \project{} beats provenance-only, proving labels alone are insufficient.
  \item \project{} beats spotlighting-style context wrapping on write-path and propagation rows.
  \item \project{} beats trust-threshold sanitizer on clean utility or attack reduction because it is
        more selective.
  \item \project{} remains competitive on latency and memory overhead.
  \item Non-RAG subsystems show real gains.
\end{itemize}

\section{Comparison Tables to Include}

\begin{longtable}{|L{0.20\textwidth}|L{0.70\textwidth}|}
\hline
\textbf{Table} & \textbf{Purpose} \\
\hline
Framework/subsystem coverage table & Show LangGraph, LlamaIndex, and CrewAI support for each memory
subsystem and whether adapter coverage is full, partial, or missing. \\
\hline
Defense comparison table & Rows: defense settings. Columns: ASR reduction, URR, CCR, BUP, p95
latency, memory overhead. Split by subsystem if space allows. \\
\hline
Prior-defense comparison table & Rows: \project{}, spotlighting, trust threshold, optional
RAGForensics, optional Task Shield. Columns: supported lifecycle stages, memory subsystems, framework
coverage, metrics improved. \\
\hline
Ablation table & Rows: full system and one-module-removed variants. Columns: ASR, CCR, propagation
depth, BUP, p95 latency. \\
\hline
Scalability table & Rows: memory size, poison ratio, top-k, agent count. Columns: slope, confidence
interval, dominant failure mode. \\
\hline
\end{longtable}

\section{Ablation Interpretation}

\begin{itemize}
  \item If removing the write gate barely changes results, the write-path contribution is weak.
  \item If removing context shaping barely changes CCR, the context-shaping contribution is weak.
  \item If removing propagation control barely changes fan-out, the distinctive multi-agent claim is
        weak.
  \item If full \project{} is only marginally better than trust-threshold sanitizer, simplify claims
        and emphasize portability and logging rather than defense quality.
  \item If clean utility drops more than 10 percent on core tasks, the paper becomes difficult to
        defend.
\end{itemize}

\chapter{Paper Skeleton}

\section{Title Candidates}

\begin{itemize}
  \item \project: A Provenance-Aware Memory Firewall for LLM Agents
  \item Defending Mutable Memory in LLM Agents with Provenance and Lineage
  \item Memory Firewalls for Stateful LLM Agents
  \item Write, Retrieve, Propagate: Securing Agent Memory Across Frameworks
\end{itemize}

\section{Abstract Draft Structure}

\begin{enumerate}
  \item State the problem: LLM agents increasingly rely on writable, persistent memory.
  \item State the gap: existing defenses are fragmented across retrieval, prompts, or tools.
  \item State the system: \project{} is a provenance-aware memory firewall for write, retrieval,
        context, tool, and propagation paths.
  \item State the method: common schema, quarantine, risk-aware retrieval, typed evidence, tool
        verification, lineage control.
  \item State the evidence: three frameworks, four memory subsystems, six scenarios, baselines,
        ablations, scalability sweeps.
  \item State the result: fill in quantitative reductions only after experiments are complete.
\end{enumerate}

\missing{} Do not write numeric claims in the abstract until final results are generated.

\section{Paper Sections}

\begin{longtable}{|L{0.15\textwidth}|L{0.25\textwidth}|L{0.25\textwidth}|L{0.25\textwidth}|}
\hline
\textbf{Section} & \textbf{What to write} & \textbf{Evidence/figures} & \textbf{Reviewer objection to
preempt} \\
\hline
Introduction & Writable memory changes agent security from one-turn prompt defense to persistent
state defense. State contributions crisply. & Motivating example and contribution list. & ``This is
just another RAG defense.'' \\
\hline
Background & Define vector/RAG, semantic, episodic, shared coordination memory. Explain public
framework support. & Subsystem table. & ``The memory taxonomy is arbitrary.'' \\
\hline
Threat Model & Define attacker capabilities, goals, assumptions, in/out scope. & Capability-to-
scenario table. & ``The attacker is too strong or vague.'' \\
\hline
System Design & Present schema, gates, scoring, context shaping, tool verification, propagation. &
Architecture figure, schema snippet, scoring equation. & ``The defense is ad hoc.'' \\
\hline
Implementation & Explain adapters and logging at high level. & Framework coverage table. &
``Framework-agnostic is not implemented.'' \\
\hline
Methodology & Describe tasks, scenarios, baselines, metrics, seeds, model parameters. & Scenario
matrix, metric table. & ``The evaluation is not reproducible.'' \\
\hline
Evaluation & Report attack reduction, contaminated context, clean utility, ablations. & Main tables
and ablation table. & ``Security comes from over-blocking.'' \\
\hline
Scalability & Show memory-size, poison-ratio, top-k, and agent-count sweeps. & Three scalability
figures. & ``It only works on toy stores.'' \\
\hline
Discussion & Analyze failures, tradeoffs, and deployment constraints. & Error analysis examples. &
``The paper hides limitations.'' \\
\hline
Related Work & Map memory architectures, memory attacks, RAG defenses, prompt injection defenses,
tool defenses, multi-agent propagation. & No figure required. & ``Novelty is unclear.'' \\
\hline
Limitations and Ethics & State synthetic setting, hosted model reproducibility, no formal guarantee,
safe stub tools. & Artifact statement. & ``The evaluation is unsafe or overclaimed.'' \\
\hline
Conclusion & Restate bounded result and why lifecycle memory defense matters. & None. & ``Conclusion
introduces new claims.'' \\
\hline
Appendix & Thresholds, prompts, scenario manifests, extra plots, adapter details. & Full artifact
details. & ``Key details omitted.'' \\
\hline
\end{longtable}

\section{Related Work Map}

\begin{itemize}
  \item Agent memory architectures: modular memory taxonomies, long-term memory, semantic/episodic
        framing, shared memory.
  \item Memory and RAG attacks: indirect prompt injection, PoisonedRAG, AgentPoison, MINJA, Zombie
        Agents, memory injection.
  \item Agent security benchmarks: InjecAgent, AgentDojo, Agent Security Bench.
  \item Prompt and retrieval defenses: spotlighting, structured queries, provenance cueing,
        trust-aware sanitization, RAGForensics.
  \item Tool and task defenses: Task Shield and task-alignment gating.
  \item Multi-agent propagation: prompt infection and remediation-style work.
\end{itemize}

\missing{} Bibliographic metadata must be verified from primary sources before paper submission.
The source report provides URLs and names, not complete citations.

\chapter{Timeline and Execution Checklist}

\section{Two-Week Timeline}

Assume work starts on May 2, 2026 and the goal is a submission-quality compact paper by May 16,
2026. This timeline is aggressive and only realistic if the MVP stays narrow.

\begin{longtable}{|L{0.11\textwidth}|L{0.32\textwidth}|L{0.27\textwidth}|L{0.20\textwidth}|}
\hline
\textbf{Date} & \textbf{Coding milestone} & \textbf{Writing/figure milestone} & \textbf{Hard rule} \\
\hline
May 2 & Freeze title, thesis, threat model, frameworks, subsystems, baselines. & Draft abstract
slots and intro claims. & No code before claims are bounded. \\
\hline
May 3 & Smoke-test LangGraph, LlamaIndex, CrewAI; define schema and logs. & Draft threat model and
related-work map. & Cut any framework that does not run cleanly after a short attempt. \\
\hline
May 4 & Implement core schema, audit logger, LangGraph adapter, write gate. & Draft architecture
subsection. & Get one defended end-to-end run. \\
\hline
May 5 & Port write path to LlamaIndex and CrewAI. & Draft framework coverage table. & Keep schema
identical across frameworks. \\
\hline
May 6 & Implement quarantine, duplicate throttling, contradiction checks. & Draft metric
definitions. & No model training. \\
\hline
May 7 & Implement retrieval reranking, source diversity, context shaping. & Draft methodology. &
Produce first retrieval table. \\
\hline
May 8 & Implement tool verifier and CrewAI propagation controller. & Draft evaluation section shell.
& Fix lineage logging before running more propagation. \\
\hline
May 9 & Run mock matrix: one model, three frameworks, two scenarios, three baselines. & Freeze
version-one claims. & Do not expand scope after this point. \\
\hline
May 10 & Add episodic scenario and flooding stress; optional AutoGen only if easy. & Start figures. &
Cut AutoGen immediately if messy. \\
\hline
May 11 & Run clean utility and conflict/recency experiments. & Fill utility and ablation tables. &
Clean utility is mandatory. \\
\hline
May 12 & Run memory-size and top-k sweeps. & Draft scalability section. & At least one scaling curve
must exist. \\
\hline
May 13 & Run agent-count propagation sweeps; optional second model if done. & Refine captions and
claims. & Do not add second model before propagation. \\
\hline
May 14 & Compute final metrics, error analysis, ablations. & Complete full paper draft. & Cut weak
modules rather than hide them. \\
\hline
May 15 & Harsh internal review, threat-model tightening, table cleanup. & Near-final submission
draft. & Make claims smaller and sharper. \\
\hline
May 16 & Reproducibility package, appendix, final compile, artifact naming. & Final submit. & No
last-minute new experiments. \\
\hline
\end{longtable}

\section{Minimum Viable Paper Path}

\begin{itemize}
  \item Three frameworks: LangGraph, LlamaIndex, CrewAI.
  \item Four subsystems: vector/RAG, semantic, episodic, shared coordination.
  \item Five defense settings: no defense, provenance-only, spotlighting, trust threshold, full
        \project{}.
  \item Six scenarios: poisoned fact, retrieved instruction, experience graft, conflict/recency,
        shared cascade, flooding.
  \item One primary model and one embedding backend.
  \item Clean utility table, main attack table, ablation table, scalability figures.
\end{itemize}

\section{Strong Version Path}

Add only after the MVP tables are complete:

\begin{itemize}
  \item AutoGen as a fourth framework.
  \item Second model transfer.
  \item RAGForensics on vector rows.
  \item Task Shield on tool rows.
  \item Letta/MemGPT archival/shared-memory case if integration is trivial.
  \item Public benchmark subset only if it plugs in cleanly on day one.
\end{itemize}

\section{Phase Checklist}

\begin{itemize}
  \item Code scaffold complete.
  \item Three framework adapters smoke-tested.
  \item Canonical schema stable.
  \item Audit logging complete.
  \item Quarantine isolation tested.
  \item Retrieval reranking logged term-by-term.
  \item Context shaping prevents raw instruction rendering.
  \item Tool verifier tested with high-impact stub tools.
  \item Propagation lineage graph generated.
  \item Baselines implemented fairly.
  \item Scenarios auto-gradable.
  \item Aggregation reproduces tables from logs.
  \item Figures generated from scripts.
\end{itemize}

\chapter{Risks, Reviewer Objections, and Mitigations}

\section{Reviewer 1: Security and System Design Reviewer}

\textbf{Major concerns.}
\begin{itemize}
  \item Threat model may be too broad if write, retrieval, tool, and propagation attacks are all
        claimed without tight boundaries.
  \item Defense may look heuristic unless each score has a clear purpose and ablations.
  \item Quarantine and trust scoring may create a false sense of security without strong audit and
        isolation guarantees.
  \item Framework adapters may be shallow wrappers that do not capture real memory flows.
\end{itemize}

\textbf{Minor concerns.}
\begin{itemize}
  \item Naming ``firewall'' may imply stronger guarantees than delivered.
  \item Tool verifier scope may be underspecified.
  \item Procedural memory is mentioned but not defended in the MVP.
\end{itemize}

\textbf{Missing evidence.}
\begin{itemize}
  \item Quarantine precision/recall.
  \item Proof that quarantined records cannot enter normal retrieval.
  \item Per-stage ablation results.
  \item Adapter coverage table.
\end{itemize}

\textbf{Weak claims.} ``Framework-agnostic,'' ``general defense,'' and ``memory firewall'' are weak
unless bounded by implementation evidence.

\textbf{Required fixes.} Freeze the threat model, show exact trust boundaries, log every decision,
and include ablations for write gate, retrieval gate, context shaping, tool verifier, and propagation.

\textbf{Borderline outcome.} The paper has a nice architecture and some vector/RAG results, but
shared-memory propagation is toy-like and clean utility is weak.

\textbf{Strong accept outcome.} The paper demonstrates real adapters, strong non-RAG results, capped
propagation under agent-count scaling, bounded overhead, and honest limitations.

\section{Reviewer 2: Machine Learning and LLM Evaluation Reviewer}

\textbf{Major concerns.}
\begin{itemize}
  \item Evaluation may rely on one model and synthetic tasks, limiting generality.
  \item Auto-graders may be brittle or biased toward the defense.
  \item Results may be sensitive to thresholds.
  \item Hosted model nondeterminism may weaken reproducibility.
\end{itemize}

\textbf{Minor concerns.}
\begin{itemize}
  \item The report does not specify model, temperature, embedding backend, or prompt templates.
  \item Three seeds may be thin if variance is high.
  \item Utility metrics may be task-specific and not comparable across frameworks.
\end{itemize}

\textbf{Missing evidence.}
\begin{itemize}
  \item Threshold sensitivity plot.
  \item Confidence intervals.
  \item Clean utility table.
  \item Second-model transfer if time allows.
  \item Grader validation examples.
\end{itemize}

\textbf{Weak claims.} Any claim that \project{} is model-general is weak without at least a second
model or a careful limitation.

\textbf{Required fixes.} Pin model parameters, log prompts and outputs, keep graders deterministic,
run threshold sweeps, and report variance.

\textbf{Borderline outcome.} The system reduces attack success but only with aggressive thresholds
that harm benign tasks.

\textbf{Strong accept outcome.} The system preserves clean utility, improvements are stable across
threshold ranges, and a second model shows the main trend persists.

\section{Reviewer 3: Conference Acceptance Reviewer}

\textbf{Major concerns.}
\begin{itemize}
  \item Novelty may be unclear in a crowded prompt-injection and RAG-defense space.
  \item The paper may be too broad for its evidence.
  \item Baseline comparison may be incomplete or unfair.
  \item The work may read as engineering middleware rather than research.
\end{itemize}

\textbf{Minor concerns.}
\begin{itemize}
  \item Related work may be under-verified because the source report gives URLs, not full citations.
  \item The title may overpromise if results are modest.
  \item Page budget may be tight for architecture, methodology, and evaluation.
\end{itemize}

\textbf{Missing evidence.}
\begin{itemize}
  \item A single table showing cross-framework and cross-subsystem results.
  \item A clear comparison with simple and strong baselines.
  \item Scalability curves that tie security outcome to scale.
  \item Error analysis explaining failure cases.
\end{itemize}

\textbf{Weak claims.} ``First,'' ``comprehensive,'' and ``general'' should be avoided unless verified
against prior work.

\textbf{Required fixes.} Make the contribution sentence brutally precise: lifecycle memory defense
across four memory subsystems and three public frameworks, evaluated under scaling pressure.

\textbf{Borderline outcome.} The paper is coherent but looks like a collection of wrappers with
small synthetic experiments.

\textbf{Strong accept outcome.} The paper changes how reviewers think about agent memory security by
showing that write, retrieve, and propagate paths must be evaluated together.

\section{Global Risk Table}

\begin{longtable}{|L{0.25\textwidth}|L{0.12\textwidth}|L{0.28\textwidth}|L{0.25\textwidth}|}
\hline
\textbf{Risk} & \textbf{Severity} & \textbf{Why serious} & \textbf{Mitigation} \\
\hline
Looks like another RAG defense & High & Most public defenses already target retrieval. & Put
semantic, episodic, and shared-memory results in the main table. \\
\hline
Threat model vague & High & Security reviewers punish undefined attacker powers. & Capability-to-
scenario table and strict out-of-scope list. \\
\hline
Defense too heuristic & High & Ad hoc filters are weak novelty. & Interpretable feature terms,
threshold sweeps, and ablations. \\
\hline
Experiments too small & High & Cross-framework claim needs real breadth. & Three frameworks, four
subsystems, multiple seeds, scaling sweeps. \\
\hline
Clean utility loss & High & Over-blocking is not a useful defense. & BUP table and security-cost
frontier. \\
\hline
Framework integration failure & Medium to high & Adapter claim collapses. & Cut optional frameworks
early; keep MVP adapters simple. \\
\hline
Scalability synthetic only & Medium to high & Reviewers may dismiss toy scaling. & Use real runtime
paths and real framework stores where possible. \\
\hline
Baseline reproduction sinkhole & Medium & Reproducing too many defenses wastes time. & Include
three clean baselines; optional prior defenses only when easy and fair. \\
\hline
\end{longtable}

\section{Final Deliverable Checklist}

Before submission, every item below must be true:

\begin{itemize}
  \item Code complete for core schema, gates, stores, adapters, baselines, scenarios, runner, metrics,
        aggregation, and plotting.
  \item LangGraph, LlamaIndex, and CrewAI adapters smoke-tested and used in main results.
  \item Experiments complete for clean, attacked, defended, and ablation conditions.
  \item Tables generated from raw logs: framework coverage, scenario matrix, main attack results,
        clean utility, ablations, baseline comparison, scalability summary.
  \item Figures generated from scripts: architecture, lineage graph, memory-size scaling,
        agent-count propagation, security-cost frontier, threshold sensitivity if space allows.
  \item Logs saved for every run: manifests, records, decisions, retrievals, contexts, tool calls,
        lineage, grader outputs, timings.
  \item Claims validated against evidence; no abstract numeric claim appears without a table or
        figure.
  \item Related work checked from primary sources and full citations completed.
  \item Threat model bounded and out-of-scope list included.
  \item Clean utility loss reported honestly.
  \item Reviewer objections addressed in introduction, methodology, evaluation, and limitations.
  \item Paper compiles from a clean checkout.
  \item Artifact README explains how to reproduce every table and figure.
  \item Comparison table included to show whether \project{} outperforms no defense, provenance-only,
        spotlighting-style wrapper, trust-threshold sanitizer, random quarantine, and optional prior
        defenses where applicable.
\end{itemize}

\section{Harsh Final Verdict}

The best version of this project is narrow, disciplined, and empirical. Build the smallest system
that undeniably shows cross-subsystem memory defense plus scaling. The fastest way to weaken the
paper is to overreach on frameworks, domains, baseline reproductions, or claims of generality. The
fastest way to make it publishable is to keep the claim bounded, the runtime pipeline crisp, the
logs complete, the baselines fair, and the scalability evidence central.

\begin{thebibliography}{99}

\bibitem{source-report}
\texttt{paper\_roadmap.pdf}, repository source report, extracted May 2, 2026.

\bibitem{coala}
Source report cited CoALA-style modular agent memory taxonomies:
\url{https://arxiv.org/abs/2309.02427}.

\bibitem{ragforensics}
Source report cited RAGForensics:
\url{https://arxiv.org/abs/2504.21668}.

\bibitem{agentpoison}
Source report cited public memory poisoning work:
\url{https://arxiv.org/abs/2402.07867}.

\bibitem{langgraph-memory}
Source report cited LangGraph/LangChain memory documentation:
\url{https://docs.langchain.com/oss/javascript/concepts/memory}.

\bibitem{indirect-prompt-injection}
Source report cited indirect prompt injection:
\url{https://arxiv.org/abs/2302.12173}.

\bibitem{prompt-infection}
Source report cited prompt infection / multi-agent propagation:
\url{https://arxiv.org/abs/2410.07283}.

\bibitem{spotlighting}
Source report cited spotlighting:
\url{https://arxiv.org/abs/2403.14720}.

\bibitem{agent-security-bench}
Source report cited agent security benchmarks:
\url{https://arxiv.org/abs/2403.02691}.

\bibitem{llamaindex-memory}
Source report cited LlamaIndex memory documentation:
\url{https://docs.llamaindex.ai/en/stable/module_guides/deploying/agents/memory/}.

\bibitem{crewai-memory}
Source report cited CrewAI memory documentation:
\url{https://docs.crewai.com/en/concepts/memory}.

\bibitem{autogen-memory}
Source report cited AutoGen memory documentation:
\url{https://microsoft.github.io/autogen/dev/user-guide/agentchat-user-guide/memory.html}.

\bibitem{memgpt}
Source report cited MemGPT/Letta memory hierarchy:
\url{https://arxiv.org/abs/2310.08560}.

\bibitem{trust-sanitizer}
Source report cited trust-aware sanitization:
\url{https://arxiv.org/abs/2601.05504}.

\end{thebibliography}

\end{document}
